\documentclass[11pt,a4paper]{article}
% Spolecny styl pro ucitelske pripravy MAT (Work/MAT)
% Verze 1.1, 2026-10-23
\NeedsTeXFormat{LaTeX2e}
\ProvidesPackage{priprava}[2026/10/23 v1.1 Sablona ucitelskych priprav MAT]

% Zakladni typografie a cestina
\RequirePackage[T1]{fontenc}
\RequirePackage[utf8]{inputenc}
\RequirePackage[czech]{babel}
\RequirePackage{lmodern}
\RequirePackage{microtype}
\RequirePackage{amsmath,amssymb}
\RequirePackage{xcolor}
\RequirePackage{enumitem}
\RequirePackage[most]{tcolorbox}
\RequirePackage{tabularx,booktabs,array}
\RequirePackage{fancyhdr}
\RequirePackage{needspace}
\RequirePackage{geometry}
\RequirePackage[hidelinks]{hyperref}

\geometry{a4paper,top=17mm,bottom=18mm,left=23mm,right=23mm}
\definecolor{notegray}{gray}{0.39}
\setlength{\parindent}{0pt}
\setlength{\parskip}{2.5pt plus 1pt minus 0.5pt}
\setlist[itemize]{topsep=2pt,itemsep=1pt,parsep=0pt,leftmargin=1.45em}
\setlist[enumerate]{topsep=2pt,itemsep=2pt,parsep=0pt,leftmargin=1.65em}
\setcounter{secnumdepth}{0}
\renewcommand{\arraystretch}{1.12}

% Zahlavi: vlevo rocnik | globalni tema, vpravo cislo: nazev hodiny
% Zapati zustava zcela prazdne.
\pagestyle{fancy}
\fancyhf{}
\renewcommand{\headrulewidth}{0.35pt}
\renewcommand{\footrulewidth}{0pt}
\setlength{\headheight}{14pt}
\newcommand{\nastavPripravu}[4]{%
  \fancyhead[L]{\small #1\enspace\textbar\enspace #2}%
  \fancyhead[R]{\small #3: #4}%
}

% Minimalisticke cernobile boxy. Kazda definice zvlast.
\tcbset{commonlesson/.style={enhanced,breakable=false,colback=white,colframe=black,
  colbacktitle=white,coltitle=black,boxrule=0.45pt,
  arc=0.8mm,outer arc=0.8mm,left=3mm,right=3mm,top=1.6mm,bottom=1.6mm,
  toptitle=1mm,bottomtitle=0.7mm,fonttitle=\bfseries,
  attach boxed title to top left={xshift=2.5mm,yshift=-1.15mm},
  boxed title style={colback=white,colframe=white,boxrule=0pt,left=1mm,right=1mm}}}
\newtcolorbox{zadaniBox}[1][]{commonlesson,title={Zadání},#1}
\newtcolorbox{prikladBox}[1][]{commonlesson,title={Příklad},#1}
\newtcolorbox{definiceBox}[1][]{commonlesson,title={Definice},#1}

% Spolecny box pro zadani a jeho reseni, bez mezery mezi castmi.
% Pouziti:
%   \begin{prikladReseniBox}
%     Zadani ulohy.
%     \tcblower
%     \begin{align*} ... \end{align*}
%   \end{prikladReseniBox}
% Spodni cast ma jemny sedy podklad a delici linku; lze v ni sazet
% vice radku matematiky. Obe casti maji jen jeden spolecny ramecek.
\tcbset{solutionlesson/.style={
  commonlesson,
  enhanced jigsaw,
  breakable,
  colbacklower=black!3,
  segmentation style={solid,draw=black!25,line width=0.35pt},
  fontlower=\small\color{notegray},
  before lower={\noindent\textbf{Řešení:}\enspace},
  leftlower=3mm,rightlower=3mm
}}
\newtcolorbox{prikladReseniBox}[1][]{solutionlesson,title={Příklad},#1}
\newtcolorbox{zadaniReseniBox}[1][]{solutionlesson,title={Zadání},#1}

% Pomucky pro autora pripravy: sedive poznamky; \reseni pro starsi pripravy
\newcommand{\poznamka}[1]{\par\begingroup\small\color{notegray}\itshape
  \textnormal{\textbf{Poznámka pro učitele.}}\enspace #1\par\endgroup}
\newcommand{\reseni}[1]{\par\begingroup\small\color{notegray}
  \textbf{Očekávané řešení:}\enspace #1\par\endgroup}
\newcommand{\blok}[1]{\needspace{5\baselineskip}\subsection*{#1}}
\newcommand{\vyrok}[1]{\uv{#1}}

\nastavPripravu{1. ročník}{Algebraické výrazy}{46}{Algebraické vzorce}
\begin{document}

% Zdroje: vlastni ucebnice Algebraicke vyrazy, Algebraicke vzorce;
% Realisticky.cz 010803 a 010804 Pouziti vzorcu pri uprave mnohoclenu.
\poznamka{Jádro 45 minut: připomenutí mocnění 5 min, odvození druhých mocnin 10 min, rozdíl čtverců 7 min, procvičení 10 min, třetí mocniny 8 min, shrnutí 5 min. Součet a rozdíl třetích mocnin uvádíme v přehledu, základní procvičení ponecháváme pro druhé a třetí mocniny dvojčlenů.}
\blok{Vzorec není trik: je to zkrácené roznásobování}
V hodině 45 jsme umocňovali opakovaným násobením. Jestliže se stejný typ výpočtu opakuje, zapíšeme jeho výsledek jako \textbf{identitu}, která platí pro všechna reálná $a,b$.
\blok{Přehled sedmi algebraických vzorců}
\begin{center}
\small
\renewcommand{\arraystretch}{1.55}
\begin{tabularx}{\linewidth}{@{}>{\raggedright\arraybackslash}p{38mm}X@{}}
\toprule
\textbf{Název} & \textbf{Vzorec}\\
\midrule
Druhá mocnina součtu & $(a+b)^2=a^2+2ab+b^2$\\
Druhá mocnina rozdílu & $(a-b)^2=a^2-2ab+b^2$\\
Rozdíl čtverců & $a^2-b^2=(a-b)(a+b)$\\
Třetí mocnina součtu & $(a+b)^3=a^3+3a^2b+3ab^2+b^3$\\
Třetí mocnina rozdílu & $(a-b)^3=a^3-3a^2b+3ab^2-b^3$\\
Součet krychlí & $a^3+b^3=(a+b)(a^2-ab+b^2)$\\
Rozdíl krychlí & $a^3-b^3=(a-b)(a^2+ab+b^2)$\\
\bottomrule
\end{tabularx}
\end{center}
\normalsize
\textbf{Všechny identity platí pro libovolná $a,b\in\mathbb R$.} Vzorce pro rozdíl čtverců a součet či rozdíl krychlí lze číst i zprava doleva při rozkládání na součin.
\begin{prikladReseniBox}[title={Příklad 1: odvoď prostřední člen}]
Vysvětli, odkud pochází člen $2ab$ ve vzorci pro $(a+b)^2$.
\tcblower
$(a+b)(a+b)=a^2+ab+ba+b^2=a^2+2ab+b^2$. Prostřední člen vzniká ze \emph{dvou} stejných součinů.
\end{prikladReseniBox}
Pro rozdíl čtverců dostaneme přímo distributivitou
\[(a-b)(a+b)=a^2+ab-ab-b^2=a^2-b^2.\]
\begin{prikladReseniBox}[title={Příklad 2: obě strany vzorce}]
Uprav $(3x-2)^2$ a rozlož $9x^2-4$ na součin.
\tcblower
$(3x-2)^2=9x^2-12x+4$; $9x^2-4=(3x-2)(3x+2)$.
\end{prikladReseniBox}
\blok{Třetí mocniny dvojčlenů}
Třetí mocniny vyplynou z násobení čtverce dvojčlenu stejným dvojčlenem. U rozdílu se střídají znaménka; součet a rozdíl krychlí naopak slouží zejména k rozkladu.
\begin{prikladReseniBox}[title={Příklad 3: odvození třetí mocniny}]
Vynásob $(a+b)^2(a+b)$ a ověř vzorec.
\tcblower
$(a^2+2ab+b^2)(a+b)=a^3+a^2b+2a^2b+2ab^2+ab^2+b^3=a^3+3a^2b+3ab^2+b^3$.
\end{prikladReseniBox}
\blok{Samostatná práce — základ}
\begin{zadaniReseniBox}[title={S1: vyber správný vzorec}]
\begin{enumerate}[label=\alph*)]
\item $(x+5)^2$;
\item $(2x-3)^2$;
\item $(4a+b)(4a-b)$;
\item $(x+2)^3$;
\item $(2y-1)^3$;
\item Rozlož $25x^2-9$.
\end{enumerate}
\tcblower
\begin{enumerate}[label=\alph*)]
\item $x^2+10x+25$;
\item $4x^2-12x+9$;
\item $16a^2-b^2$;
\item $x^3+6x^2+12x+8$;
\item $8y^3-12y^2+6y-1$;
\item $(5x-3)(5x+3)$.
\end{enumerate}
\end{zadaniReseniBox}
\poznamka{Nech studenty obhájit výběr vzorce. V každém případě rozlišuj výraz typu $(a-b)^2$ od rozdílu $a^2-b^2$. Proti časté chybě stačí protipříklad $a=2,b=1$.}
\begin{prikladReseniBox}[title={Příklad 4: součet a rozdíl krychlí}]
Rozlož $x^3+8$ a $27a^3-b^3$ na součin.
\tcblower
$x^3+8=x^3+2^3=(x+2)(x^2-2x+4)$;\\
$27a^3-b^3=(3a)^3-b^3=(3a-b)(9a^2+3ab+b^2)$.
\end{prikladReseniBox}
\blok{Volitelná rezerva}
\begin{zadaniReseniBox}[title={R1: poznáš skrytý čtverec?}]
Zapiš $x^2-6x+9$ jako druhou mocninu dvojčlenu a ověř roznásobením.
\tcblower
$x^2-6x+9=(x-3)^2$. Po roznásobení opět získáme původní trojčlen.
\end{zadaniReseniBox}
\blok{Domácí procvičení}
\begin{zadaniReseniBox}[title={D1}]
\begin{enumerate}[label=\alph*)]
\item $(3x+2)^2$;
\item $(a-2b)^2$;
\item $(x+3)(x-3)$;
\item $(x-2)^3$;
\item $(2a+b)^3$;
\item Rozlož $x^3-27$;
\item Rozlož $8a^3+b^3$.
\end{enumerate}
\tcblower
\begin{enumerate}[label=\alph*)]
\item $9x^2+12x+4$;
\item $a^2-4ab+4b^2$;
\item $x^2-9$;
\item $x^3-6x^2+12x-8$;
\item $8a^3+12a^2b+6ab^2+b^3$;
\item $(x-3)(x^2+3x+9)$;
\item $(2a+b)(4a^2-2ab+b^2)$.
\end{enumerate}
\end{zadaniReseniBox}
\end{document}
